% Abstract - adapted for tufte class
\chapter*{Abstract}
\addcontentsline{toc}{chapter}{Abstract}

Individual nuclear spins in the solid state are attractive candidates for quantum computing and high-precision spectroscopy thanks to their long coherence times. So far, the only mechanisms to detect individual nuclear spins is using an electron spin in close proximity, which couples via hyperfine interaction and acts as a local antenna. Therefore, individual nuclear spin detection requires single electron spin readout.

Until recently, detection of single electron spins has been limited to a handful of systems with specific readout mechanisms, which has limited the range of study. However, microwave photon counting provides a more general approach to single electron spin detection. By coupling the electron spin to a low-mode-volume, high-quality-factor resonator, the radiative relaxation rate of the spin is enhanced by the Purcell effect. The resulting fluorescence is detected by a single microwave photon detector, a high sensitivity device based on transmon qubit. 

Using microwave fluorescence detection we identify single \W nuclear spins that are coupled to individual \Er electron spin in \Ca. We develop methods for single-shot readout of the nuclear spin state, deterministic preparation using solid-effect dynamic nuclear polarization and hyperfine spectroscopy. Using the measured values of the hyperfine coupling we are able to tentatively assign the location of the nuclear spin with respect to the \Er spin.

In addition, we use microwave stimulated Raman gates to create single- and two-qubit gates in a two \W nuclear spin quantum register. By using a Raman process, we do not significantly populate the electron spin excited state, which enables coherent control of the nuclear spin state while preserving its long decoherence time of $\sim1$~s. These techniques demonstrate readout, preparation and control of nuclear spin quantum register using only microwave signals.

Finally, we detect an unexpected nuclear spin-9/2 coupled to an \Er spin. Through precision NMR spectroscopy we chemically identify it as a \Nb nuclear spin and fully reconstruct its quadrupole tensor, which we use to determine its substitution site. Furthermore, we identify a previously unobserved coupling term between the electron spin and the nuclear quadrupole, and we set a tighter bound on the hexadecapolar interaction.

Overall, the methods presented in this thesis demonstrate the capabilities of microwave photon counting for quantum computing and precision spectroscopy.